\documentclass[book.tex]{subfiles}

\begin{document}
\chapter{A Simple Detailed QED Calculation}
\label{chap:qed_calculation}
\noindent\rule{11cm}{0.4pt} \\
\textbf{Prerequisites:} \autoref{chap:qft_basics}  \\
\textbf{Difficulty Level:} ** \\
\noindent\rule{11cm}{0.4pt}
\vspace{5mm}

%\textcolor{red}{TODO: Cite Peskin and Schroeder}

In this chapter we study the scattering reaction
\begin{align}
    e^+e^- \mapsto \mu^+\mu^-
\end{align}
%This process is represented by the diagram below in Figure~\ref{fig:qed_basic}.
% The feynman diagrams are generated in separate overleaf feynman_diagrams due to compilation issues for now. TODO: Unify style files

\begin{figure}
    \centering
    % \includegraphics[width=0.5\textwidth]{figures/Differentiable Physics/QFT ML/QED Calculation/electronmuon2.png}
    \includegraphics[width=0.5\textwidth]{figures/Differentiable Physics/QFT ML/QED Calculation/electronmuon.png}
    \caption{$e^+e^-\mapsto \mu^+\mu^-$ scattering process.}
    \label{fig:qed_basic}
\end{figure}

This diagram has the following amplitude %(\textcolor{red}{TODO: Write down feynman diagram rules}

\begin{align}
    \mathcal{M} &= \bar{v}^{s'}(p')(-ie\gamma^{\mu})u^s(p)\left ( \frac{-ig_{\mu\nu}}{q^2}\right ) \bar{u}^r(k)(-ie\gamma^\nu)v^{r'}(k')
\end{align}

\section{Real Space Feynman Rules}

Terms are computed from Feynman diagrams according to the Feynman rules. We list the Feynman rules for QED

\begin{itemize}
    \item Integration coordinates $x_j$ are represented by vertices
    \item Bosonic propagators are represented by wiggly lines connecting two vertices.
    \item Fermionic propagators are represented by straight lines connecting two vertices.
    \item A bosonic field $A_\mu(x_i)$ is given by a wiggly line attached to a point.
    \item A fermionic field $\psi(x_i)$ is given by a straight line attached to a point with the arrow towards the point. 
    \item An anti-fermionic field $\overline{\psi}(x_i)$ is given by a straight line attached to the point with the arrow pointed away from the point.
\end{itemize}

\section{Momentum Space Feynman Rules}

It is often useful to compute Feynman rules in momentum space rather than in real space. Analogous transformations of these rules exist for momentum space.

\end{document}
